\documentclass[aspectratio=169]{beamer}

\usetheme{Madrid}
\usepackage[spanish]{babel}
\usepackage[utf8]{inputenc}
\usepackage{textcomp}
\usepackage[T1]{fontenc}
\usepackage{lmodern}
\usepackage{graphicx}
\usepackage{booktabs}
\usepackage{verbatim}
\usepackage{amssymb}
\usepackage{pifont}
\usepackage{moreverb}
\usepackage[most]{tcolorbox}
\usepackage{tikz}
\usetikzlibrary{shapes, arrows, positioning}

% Configuración de TikZ
\tikzstyle{block} = [
  rectangle, draw, rounded corners,
  text centered, minimum height=1cm,
  minimum width=2.5cm, font=\small, fill=gray!5
]
\tikzstyle{arrow} = [->, thick]

% Configuración de tcolorbox
\tcbset{
  colback=gray!5, colframe=gray!50,
  boxrule=0.4pt, sharp corners,
}

\title[BookMate]{BookMate — Sistema de Recomendación de Libros Académicos}
\subtitle{Diseño Detallado: Clases, Secuencias, Componentes y Despliegue}
\author[GRUPO 6]{Delgado R., G. \and Osorio M., A. \and Rojas A., J. \and Torres R., J. \and Valverde G., Y. \and Villanueva A., F.}
\institute[FC-UNI]{UNIVERSIDAD NACIONAL DE INGENIERÍA}
\titlegraphic{\includegraphics[height=1.4cm]{example-image}}
% NOTA: Reemplazar 'example-image' con 'Imagenes/logo_uni.png' cuando esté disponible

\date{Noviembre 2025}

\begin{document}

% Portada
\begin{frame}
\titlepage
\end{frame}

% Agenda
\begin{frame}
\frametitle{Agenda}
\tableofcontents
\end{frame}

\section{Modelo de Diseño}

\begin{frame}
\frametitle{Transformación: Análisis → Diseño}

\begin{columns}
\begin{column}{0.5\textwidth}
\textbf{Análisis (S11):}
\begin{itemize}
\item Entidades conceptuales
\item Controladores conceptuales
\item Sin tecnologías específicas
\item Arquitectura conceptual
\end{itemize}
\end{column}

\begin{column}{0.5\textwidth}
\textbf{Diseño (S15):}
\begin{itemize}
\item Clases JPA con anotaciones
\item Services y Controllers reales
\item Repositories Spring Data
\item DTOs para transferencia
\item Métodos con firmas completas
\end{itemize}
\end{column}
\end{columns}

\vspace{0.5cm}

\begin{center}
\begin{tikzpicture}[node distance=3cm]
\node[block, fill=LightBlue] (analisis) {Análisis\\Conceptual};
\node[block, fill=LightGreen, right=of analisis] (diseno) {Diseño\\Detallado};
\draw[arrow] (analisis) -- node[above] {Transformación} (diseno);
\end{tikzpicture}
\end{center}

\end{frame}

\begin{frame}
\frametitle{Mapeo: Entidades Conceptuales → Clases JPA}

\begin{table}
\centering
\tiny
\begin{tabular}{p{3cm}p{4cm}p{4cm}}
\toprule
\textbf{Entidad Conceptual} & \textbf{Clase JPA} & \textbf{Anotaciones Principales} \\
\midrule
Libro & \texttt{Libro} & \texttt{@Entity}, \texttt{@ManyToOne}, \texttt{@OneToOne} \\
Autor & \texttt{Autor} & \texttt{@Entity}, \texttt{@OneToMany} \\
Representación Semántica & \texttt{Embedding} & \texttt{@Entity}, \texttt{@OneToOne} \\
Usuario & \texttt{Usuario} & \texttt{@Entity}, \texttt{@Enumerated} \\
\bottomrule
\end{tabular}
\end{table}

\vspace{0.3cm}

\textbf{Ejemplo: Libro}
\begin{verbatim}
@Entity
@Table(name = "libros")
public class Libro {
    @Id
    @GeneratedValue(strategy = IDENTITY)
    private Long id;
    
    @ManyToOne(fetch = LAZY)
    private Autor autor;
    
    @OneToOne(mappedBy = "libro")
    private Embedding embedding;
}
\end{verbatim}

\end{frame}

\begin{frame}
\frametitle{Mapeo: Controladores → Services/Controllers}

\begin{table}
\centering
\tiny
\begin{tabular}{p{3.5cm}p{3.5cm}p{3.5cm}}
\toprule
\textbf{Controlador Conceptual} & \textbf{Service} & \textbf{Controller} \\
\midrule
Gestor de Catálogo & \texttt{LibrosService} & \texttt{LibrosController} \\
Generador de Recomendaciones & \texttt{RecommendationService} & \texttt{RecommendationController} \\
Analizador Semántico & \texttt{AIServiceClient} & (Cliente HTTP) \\
Calculador Heurístico & \texttt{HeuristicRecommender} & (Service) \\
\bottomrule
\end{tabular}
\end{table}

\vspace{0.3cm}

\textbf{Separación de responsabilidades:}
\begin{itemize}
\item \textbf{Controller:} Maneja HTTP, validación, respuestas
\item \textbf{Service:} Lógica de negocio, coordinación
\item \textbf{Repository:} Acceso a datos
\end{itemize}

\end{frame}

\section{Diagrama de Clases de Diseño}

\begin{frame}
\frametitle{Diagrama de Clases Completo}

\begin{center}
\includegraphics[width=0.9\textwidth]{example-image}
% NOTA: Reemplazar con diagrama de clases de diseño completo
% Mostrar:
% - Entidades JPA (Libro, Autor, Embedding, Usuario)
% - Repositories (interfaces Spring Data)
% - Services (LibrosService, RecommendationService, etc.)
% - Controllers (LibrosController, RecommendationController, etc.)
% - DTOs (LibroDTO, AutorDTO, RecomendacionDTO)
% - Relaciones entre clases
% - Anotaciones JPA y Spring
\end{center}

\textbf{Clases principales:}
\begin{itemize}
\item \textbf{Entidades:} Libro, Autor, Embedding, Usuario
\item \textbf{Repositories:} LibrosRepository, AutorRepository, EmbeddingRepository
\item \textbf{Services:} LibrosService, RecommendationService, AIServiceClient, HeuristicRecommender
\item \textbf{Controllers:} LibrosController, AutorController, RecommendationController
\item \textbf{DTOs:} LibroDTO, AutorDTO, RecomendacionDTO
\end{itemize}

\end{frame}

\begin{frame}
\frametitle{Relaciones Entre Clases}

\textbf{Relaciones principales:}

\begin{itemize}
\item \textbf{Libro → Autor:} \texttt{@ManyToOne} (muchos libros por autor)
\item \textbf{Autor → Libro:} \texttt{@OneToMany} (un autor tiene muchos libros)
\item \textbf{Libro → Embedding:} \texttt{@OneToOne} (un embedding por libro)
\item \textbf{Controller → Service:} Inyección de dependencias
\item \textbf{Service → Repository:} Inyección de dependencias
\end{itemize}

\vspace{0.3cm}

\begin{center}
\includegraphics[width=0.7\textwidth]{example-image-a}
% NOTA: Reemplazar con diagrama de relaciones entre clases
% Mostrar relaciones ManyToOne, OneToMany, OneToOne
\end{center}

\end{frame}

\section{Diagrama de Secuencia de Diseño}

\begin{frame}
\frametitle{Secuencia: UC-06 Obtener Recomendaciones (IA)}

\begin{center}
\includegraphics[width=0.85\textwidth]{example-image-b}
% NOTA: Reemplazar con diagrama de secuencia de diseño de UC-06
% Mostrar flujo completo:
% Usuario → Frontend → RecommendationController → RecommendationService
% → AIServiceClient → Servicio IA → AIServiceClient → RecommendationService
% → LibrosRepository → PostgreSQL → RecommendationService → Controller → Frontend → Usuario
\end{center}

\textbf{Flujo principal:}
\begin{enumerate}
\item Usuario solicita recomendaciones
\item Controller invoca Service
\item Service verifica disponibilidad de IA
\item Service solicita similitud a IA
\item Service consulta libros en BD
\item Service transforma a DTOs
\item Controller retorna JSON
\end{enumerate}

\end{frame}

\begin{frame}
\frametitle{Secuencia: Fallback a Heurísticas}

\begin{center}
\includegraphics[width=0.8\textwidth]{example-image-c}
% NOTA: Reemplazar con diagrama de secuencia del flujo alternativo
% Mostrar:
% RecommendationService → AIServiceClient (verificar) → false
% → HeuristicRecommender → LibrosRepository → cálculo de puntuaciones
% → RecommendationService → Controller → Frontend
\end{center}

\textbf{Flujo alternativo:}
\begin{enumerate}
\item Service verifica IA → no disponible
\item Service activa HeuristicRecommender
\item HeuristicRecommender calcula puntuaciones
\item HeuristicRecommender ordena y selecciona top 6
\item Service retorna recomendaciones heurísticas
\end{enumerate}

\vspace{0.3cm}

\begin{tcolorbox}[colback=green!5, colframe=green!50]
\textbf{Ventaja:} Sistema siempre retorna recomendaciones, garantizando alta disponibilidad.
\end{tcolorbox}

\end{frame}

\section{Diagrama de Colaboración}

\begin{frame}
\frametitle{Colaboración: UC-06}

\begin{center}
\includegraphics[width=0.8\textwidth]{example-image}
% NOTA: Reemplazar con diagrama de colaboración de UC-06
% Mostrar objetos y mensajes entre ellos:
% recController → recService → aiClient → librosRepo → libro123, libro456, etc.
\end{center}

\textbf{Objetos participantes:}
\begin{itemize}
\item \textbf{recController:} RecommendationController
\item \textbf{recService:} RecommendationService
\item \textbf{aiClient:} AIServiceClient
\item \textbf{heuristic:} HeuristicRecommender
\item \textbf{librosRepo:} LibrosRepository
\item \textbf{libro123, libro456, etc.:} Instancias de Libro
\end{itemize}

\textbf{Enfoque:}
- Interacción entre objetos en tiempo de ejecución
- Mensajes entre objetos
- Colaboración para cumplir caso de uso

\end{frame}

\section{Diagrama de Componentes}

\begin{frame}
\frametitle{Componentes del Sistema}

\begin{center}
\includegraphics[width=0.85\textwidth]{example-image-a}
% NOTA: Reemplazar con diagrama de componentes UML
% Mostrar:
% - Frontend (HTML/CSS/JS, Bootstrap)
% - Backend Spring Boot (Controllers, Services, Repositories, Entities, DTOs, Security)
% - Base de Datos (PostgreSQL)
% - Servicio IA (Flask App, Sentence Transformers, Embeddings Module, Similarity Module)
% - Interfaces (HTTP REST, JDBC, JPA)
\end{center}

\textbf{Componentes principales:}
\begin{itemize}
\item \textbf{Frontend:} HTML/CSS/JS, Bootstrap 5.1.3
\item \textbf{Backend:} Controllers, Services, Repositories, Entities, DTOs, Security
\item \textbf{Base de Datos:} PostgreSQL 16
\item \textbf{Servicio IA:} Flask App, Sentence Transformers, Embeddings, Similarity
\end{itemize}

\end{frame}

\section{Diagrama de Despliegue}

\begin{frame}
\frametitle{Arquitectura de Despliegue}

\begin{center}
\includegraphics[width=0.9\textwidth]{example-image-b}
% NOTA: Reemplazar con diagrama de despliegue completo
% Mostrar:
% - Docker Host
% - Red: bookmate-network
% - Contenedor: bookmate-backend (puerto 8080)
% - Contenedor: postgresql (puerto 5432, volumen postgres-data)
% - Contenedor: bookmate-ai-service (puerto 5000)
% - Usuario (navegador) conectándose al backend
\end{center}

\textbf{Contenedores Docker:}
\begin{itemize}
\item \textbf{bookmate-backend:} Spring Boot (puerto 8080)
\item \textbf{postgresql:} PostgreSQL 16 (puerto 5432, volumen persistente)
\item \textbf{bookmate-ai-service:} Python Flask (puerto 5000)
\end{itemize}

\textbf{Orquestación:}
- Docker Compose
- Red: bookmate-network
- Volúmenes: postgres-data

\end{frame}

\begin{frame}
\frametitle{Configuración de Despliegue}

\textbf{Docker Compose:}
\begin{verbatim}
services:
  backend:
    image: bookmate-backend:1.0.0
    ports: ["8080:8080"]
    depends_on: [postgres, ai-service]
  
  postgres:
    image: postgres:16
    ports: ["5432:5432"]
    volumes: [postgres-data:/var/lib/postgresql/data]
  
  ai-service:
    image: bookmate-ai-service:1.0.0
    ports: ["5000:5000"]
\end{verbatim}

\textbf{Características:}
\begin{itemize}
\item Red compartida: bookmate-network
\item Volumen persistente para PostgreSQL
\item Health checks configurados
\item Variables de entorno para configuración
\end{itemize}

\end{frame}

\section{Diagramas de Estados}

\begin{frame}
\frametitle{Estados: Ciclo de Vida JPA (Libro)}

\begin{center}
\includegraphics[width=0.75\textwidth]{example-image-c}
% NOTA: Reemplazar con diagrama de estados del ciclo de vida JPA
% Mostrar:
% Transient → Managed (save/persist)
% Managed → Detached (clear/close)
% Detached → Managed (merge)
% Managed → Removed (remove)
% Removed → [*] (commit)
\end{center}

\textbf{Estados JPA:}
\begin{itemize}
\item \textbf{Transient:} Objeto nuevo, no persistido
\item \textbf{Managed:} Asociado a EntityManager, sincronizado con BD
\item \textbf{Detached:} Desasociado pero existe en BD
\item \textbf{Removed:} Marcado para eliminación
\end{itemize}

\end{frame}

\begin{frame}
\frametitle{Estados: Embedding}

\begin{center}
\includegraphics[width=0.7\textwidth]{example-image}
% NOTA: Reemplazar con diagrama de estados de Embedding
% Mostrar:
% Pendiente → Generando → Almacenado
% Almacenado → Actualizando → Almacenado
% Generando → Error → Generando (reintento)
\end{center}

\textbf{Estados:}
\begin{itemize}
\item \textbf{Pendiente:} Libro creado, esperando generación
\item \textbf{Generando:} Procesando sinopsis con Sentence Transformers
\item \textbf{Almacenado:} Embedding disponible en BD
\item \textbf{Actualizando:} Regenerando (sinopsis cambió)
\item \textbf{Error:} Fallo en generación (reintento programado)
\end{itemize}

\end{frame}

\begin{frame}
\frametitle{Estados: Solicitud de Recomendación}

\begin{center}
\includegraphics[width=0.7\textwidth]{example-image-a}
% NOTA: Reemplazar con diagrama de estados de solicitud de recomendación
% Mostrar:
% Solicitada → VerificandoIA → ProcesandoConIA/ProcesandoHeuristico
% → Completada → Entregada
\end{center}

\textbf{Estados:}
\begin{enumerate}
\item \textbf{Solicitada:} Petición recibida
\item \textbf{VerificandoIA:} Health check del servicio IA
\item \textbf{ProcesandoConIA:} Cálculo de similitud semántica
\item \textbf{ProcesandoHeuristico:} Cálculo de puntuaciones (fallback)
\item \textbf{Completada:} Recomendaciones generadas
\item \textbf{Entregada:} Respuesta enviada al usuario
\end{enumerate}

\end{frame}

\section{Patrones de Diseño Aplicados}

\begin{frame}
\frametitle{Resumen de Patrones}

\textbf{Patrones implementados:}

\begin{enumerate}
\item \textbf{Strategy Pattern:}
   \begin{itemize}
   \item Recomendaciones: IA vs Heurísticas
   \item Permite cambiar algoritmo en tiempo de ejecución
   \end{itemize}

\item \textbf{Repository Pattern:}
   \begin{itemize}
   \item Abstracción de acceso a datos
   \item Spring Data JPA implementa automáticamente
   \end{itemize}

\item \textbf{Service Layer Pattern:}
   \begin{itemize}
   \item Separación de lógica de negocio
   \item Facilita testing y reutilización
   \end{itemize}

\item \textbf{DTO Pattern:}
   \begin{itemize}
   \item Transferencia de datos entre capas
   \item Control de qué datos se exponen
   \end{itemize}

\item \textbf{Dependency Injection:}
   \begin{itemize}
   \item Spring inyecta dependencias automáticamente
   \item Reduce acoplamiento
   \end{itemize}
\end{enumerate}

\end{frame}

\section{Resumen y Conclusiones}

\begin{frame}
\frametitle{Resumen del Diseño Detallado}

\begin{columns}
\begin{column}{0.5\textwidth}
\textbf{Clases de Diseño:}
\begin{itemize}
\item 4 entidades JPA
\item 3 repositories
\item 6 services
\item 3 controllers
\item 3+ DTOs
\end{itemize}
\end{column}

\begin{column}{0.5\textwidth}
\textbf{Diagramas Generados:}
\begin{itemize}
\item Diagrama de clases
\item Diagrama de secuencia
\item Diagrama de colaboración
\item Diagrama de componentes
\item Diagrama de despliegue
\item Diagramas de estados
\end{itemize}
\end{column}
\end{columns}

\vspace{0.5cm}

\begin{tcolorbox}[colback=blue!5, colframe=blue!50, title=Objetivo Cumplido]
Hemos transformado el modelo de análisis conceptual en un diseño de implementación completo con clases reales, métodos específicos y arquitectura física detallada. El sistema está listo para implementación.
\end{tcolorbox}

\end{frame}

\begin{frame}
\frametitle{Logros del Diseño Detallado}

\begin{enumerate}
\item ✅ \textbf{Modelo de diseño} completo con transformación análisis→diseño
\item ✅ \textbf{Diagrama de clases} con clases Java reales y relaciones
\item ✅ \textbf{Diagrama de secuencia} con llamadas reales y flujos alternativos
\item ✅ \textbf{Diagrama de colaboración} mostrando interacción entre objetos
\item ✅ \textbf{Diagrama de componentes} con tecnologías específicas
\item ✅ \textbf{Diagrama de despliegue} con Docker Compose
\item ✅ \textbf{Diagramas de estados} para entidades JPA y procesos
\item ✅ \textbf{Patrones de diseño} aplicados y justificados
\end{enumerate}

\vspace{0.5cm}

\begin{tcolorbox}[colback=green!5, colframe=green!50, title=Próximos Pasos]
Semana 16: Implementación final, pruebas, despliegue y demo completa del sistema.
\end{tcolorbox}

\end{frame}

\begin{frame}
\frametitle{¡Preguntas?}

\begin{center}
\Large ¡Gracias por su atención! \\

\vspace{1cm}

\textbf{¿Preguntas sobre el diseño detallado?}

\vspace{1cm}

\normalsize
\textbf{Grupo 6}\\
Universidad Nacional de Ingeniería\\
CC341 - Ingeniería de Software\\
Ciclo 2025-2

\vspace{0.5cm}

\textbf{Próxima entrega:} Semana 16 - Examen Final (Demo, Planes, Sistema Completo)

\end{center}

\end{frame}

\end{document}

